\subsection{From the Fibonacci presentation of scale-counting configurations to the natural numbers (quotients and invariants)}\label{subsec:emergent-arithmetic-fib-congruence}

The main point of this subsection is that the arithmetic structure is not postulated in advance. Rather, it is forced by taking the quotient of the free scale-counting configuration space by the local Fibonacci relations, or equivalently by passing to the corresponding invariant. Zeckendorf normal form is then a computable system of canonical representatives for the quotient.

Let
\[
\mathcal{D}:=\Bigl\{a=(a_1,a_2,\dots): a_k\in\NN,\ \text{only finitely many }a_k\neq 0\Bigr\}.
\]
We regard \(a\in\mathcal{D}\) as a counting configuration recording \(a_k\) occurrences at scale \(k\). Equipped with pointwise addition, \(\mathcal{D}\) is the free commutative monoid on the generators \(e_k\), where \(e_k\) denotes the \(k\)-th standard basis vector.

\begin{definition}[Congruence generated by the Fibonacci relations]\label{def:fib-congruence}
Let \(\equiv\) be the smallest congruence on the commutative monoid \(\mathcal{D}\) such that
\begin{enumerate}
  \item \textbf{(base congruence)} \(e_2\equiv 2e_1\), which folds repeated occurrences at the same weight into a single occurrence at the next scale;
  \item \textbf{(recursive congruence)} for every \(k\ge 1\),
\[
e_{k+2}\equiv e_{k+1}+e_k.
\]
\end{enumerate}
Equivalently, \(\equiv\) is generated at the level of scale-counting configurations by the base fold together with the Fibonacci recursion, and is then closed under addition.
\leanverified{paper\_val\_base\_e2}
\leanverified{paper\_val\_recursion\_e}
\end{definition}

\begin{definition}[Value homomorphism]\label{def:val-on-D}
Define \(\mathrm{Val}:\mathcal{D}\to\NN\) by
\[
\mathrm{Val}(a):=\sum_{k\ge 1} a_k F_{k+1}.
\]
This is a commutative-monoid homomorphism, and it is compatible with the earlier definition of \(\mathrm{Val}\) on \(X_\infty\) when one views an admissible sequence \(c\in X_\infty\) as the special case \(a_k=c_k\).
\leanpartial{valE}{Formalized only for single basis vectors \(e_k\); the general definition of \(\mathrm{Val}\) on arbitrary elements of \(\mathcal{D}\) remains to be completed}
\end{definition}

\begin{proposition}[Value invariance of Fibonacci congruence]\label{prop:val-invariant}
If \(a\equiv b\), then \(\mathrm{Val}(a)=\mathrm{Val}(b)\). Consequently, \(\mathrm{Val}\) induces a well-defined homomorphism on the quotient monoid
\[
\overline{\mathrm{Val}}:\mathcal{D}/\!\equiv\ \to\ \NN.
\]
\end{proposition}

\begin{proof}
It suffices to verify that the generating relations preserve value. For the recursive relation \(e_{k+2}\equiv e_{k+1}+e_k\),
\[
\mathrm{Val}(e_{k+2})=F_{k+3}=F_{k+2}+F_{k+1}=\mathrm{Val}(e_{k+1}+e_k).
\]
For the base relation \(e_2\equiv 2e_1\),
\[
\mathrm{Val}(e_2)=F_3=2=2F_2=\mathrm{Val}(2e_1).
\]
The general claim follows because \(\equiv\) is closed under addition and \(\mathrm{Val}\) is additive.
\end{proof}

\begin{theorem}[Emergence of the natural numbers: isomorphism of the quotient monoid]\label{thm:monoid-quotient-is-N}
The map \(\overline{\mathrm{Val}}:\mathcal{D}/\!\equiv\to\NN\) is an isomorphism of commutative monoids. Moreover, \(X_\infty\) furnishes a unique canonical representative for every congruence class: for each \(a\in\mathcal{D}\), there exists a unique \(c\in X_\infty\) such that \(a\equiv c\) and \(\mathrm{Val}(a)=\mathrm{Val}(c)\).
\leanpartial{paper\_val\_e\_index\_eq\_fib\_succ}{The value-side identity \([e_k]=F_{k+1}\varepsilon\) has been formalized; the quotient isomorphism and uniqueness of canonical representatives remain to be formalized}
\leanverified{valE\_zero\_one\_two\_table}
\leanverified{valE\_pos\_of\_ge\_one}
\leanverified{valE\_mono}
\leanverified{valE\_strict\_mono\_of\_ge\_two}
\leanverified{paper\_valE\_milestone\_package}
\end{theorem}

\begin{proof}
\begin{itemize}
  \item \textbf{Surjectivity.} For any \(N\in\NN\), let \(Z(N)=(c_1,c_2,\dots)\in X_\infty\) be its Zeckendorf expansion \cite{Zeckendorf1972}. Then \(\mathrm{Val}(Z(N))=N\), so \(\overline{\mathrm{Val}}\) is surjective.
  \item \textbf{Identification of all classes with a single generator.} In the quotient monoid \(\mathcal{D}/\!\equiv\), let \(\varepsilon:=[e_1]\). The base congruence gives \([e_2]=2\varepsilon\), and the recursive congruence yields by induction
\[
[e_{k+2}]=[e_{k+1}]+[e_k]
\qquad\Longrightarrow\qquad
[e_k]=F_{k+1}\varepsilon
\]
for every \(k\ge 1\). Hence for any \(a=\sum_{k\ge 1} a_k e_k\in\mathcal{D}\),
\[
[a]=\sum_{k\ge 1} a_k [e_k]
=\left(\sum_{k\ge 1} a_k F_{k+1}\right)\varepsilon
=\mathrm{Val}(a)\,\varepsilon.
\]
Thus each congruence class is completely determined by its value, and \(\overline{\mathrm{Val}}([a])=\mathrm{Val}(a)\) is precisely the monoid isomorphism sending \(n\varepsilon\) to \(n\).
  \item \textbf{Canonical representatives.} Given \(a\in\mathcal{D}\), let \(N=\mathrm{Val}(a)\) and set \(c=Z(N)\in X_\infty\). The previous step gives
\[
[a]=N\varepsilon=[c],
\]
so \(a\equiv c\). Uniqueness follows from uniqueness of Zeckendorf expansion: if \(c,c'\in X_\infty\) and \(c\equiv c'\), then \(\mathrm{Val}(c)=\mathrm{Val}(c')\), hence \(c=c'\).
\end{itemize}
\end{proof}

\begin{definition}[Stable folding map]\label{def:stable-fold}
For \(a\in\mathcal{D}\), define
\[
\Fold_\infty(a):=Z\bigl(\mathrm{Val}(a)\bigr)\in X_\infty.
\]
Equivalently, \(\Fold_\infty\) is the canonical-section map obtained by taking the value invariant of a finite counting configuration and then returning its Zeckendorf normal representative.
\end{definition}

\begin{proposition}[Finite local reachability of the Zeckendorf fold]\label{prop:finite-local-fold-normalization}
For every \(a\in\mathcal{D}\), the normal form \(\Fold_\infty(a)
=Z(\mathrm{Val}(a))\) is reachable by a finite sequence of value-preserving local Fibonacci rewrites. More precisely, allow the elementary moves, in either direction and in any additive context,
\[
e_2\longleftrightarrow 2e_1,
\qquad
e_{k+2}\longleftrightarrow e_{k+1}+e_k\quad(k\ge 1).
\]
Then there is an explicit terminating reachability procedure which starts at \(a\), ends at \(\Fold_\infty(a)\), and every intermediate configuration has the same value as \(a\). This is a finite local realization of the canonical section of the quotient; it is not asserted here to be a confluent orientation of the symmetric rewrite system.
\end{proposition}

\begin{proof}
The procedure has two finite phases. First expand every occurrence of every basis vector down to the first scale. Namely, replace each occurrence of \(e_2\) by \(2e_1\), and for \(j\ge 3\) replace each occurrence of \(e_j\) by \(e_{j-1}+e_{j-2}\), continuing recursively until no basis vector other than \(e_1\) remains. This phase terminates because each replacement strictly lowers the largest index appearing in the occurrence being expanded, and the input has finite support with finite multiplicities. The result is exactly
\[
\mathrm{Val}(a)e_1,
\]
since the same recursion gives \(e_j\rightsquigarrow F_{j+1}e_1\) for each \(j\), with the base case \(e_2\rightsquigarrow 2e_1=F_3e_1\).

Second compute the Zeckendorf expansion \(c=Z(\mathrm{Val}(a))\) by the greedy algorithm. Because \(c\) has finite support, expanding each nonzero basis vector in \(c\) by the same downward recursion gives a finite value-preserving rewrite sequence
\[
c\rightsquigarrow \mathrm{Val}(c)e_1=\mathrm{Val}(a)e_1.
\]
Reverse this finite sequence. Concatenating the first phase with the reversed second phase gives a finite sequence of the stated elementary local moves from \(a\) to \(c=\Fold_\infty(a)\).

Each elementary move preserves value: this is exactly the verification in Proposition \ref{prop:val-invariant}. Hence every intermediate configuration has value \(\mathrm{Val}(a)\). The output is admissible and unique by Theorem \ref{thm:zeckendorf-unique}; therefore the procedure computes the stated fold map and terminates.
\end{proof}

\begin{theorem}[Classification of symmetric local rewrite orbits]\label{thm:symmetric-rewrite-orbits-value-fibers}
For \(a,b\in\mathcal{D}\), the following are equivalent:
\begin{enumerate}[label=\textup{(\roman*)}]
  \item \(a\) and \(b\) are connected by a finite sequence of the symmetric local moves
\[
e_2\longleftrightarrow 2e_1,
\qquad
e_{k+2}\longleftrightarrow e_{k+1}+e_k\quad(k\ge 1),
\]
applied in additive contexts;
  \item \(a\equiv b\) in the congruence of Definition \ref{def:fib-congruence};
  \item \(\mathrm{Val}(a)=\mathrm{Val}(b)\).
\end{enumerate}
Consequently every symmetric rewrite orbit contains exactly one Zeckendorf word, namely \(\Fold_\infty(a)\). This is the canonical-section statement used in the sequel; no terminating confluent orientation of the symmetric system is being assumed.
\end{theorem}

\begin{proof}
The equivalence of (i) and (ii) is the definition of the congruence generated by the displayed relations: closing the elementary local moves under addition and finite concatenation gives exactly the least congruence containing the Fibonacci relations.

If (ii) holds, Proposition \ref{prop:val-invariant} gives \(\mathrm{Val}(a)=\mathrm{Val}(b)\), so (iii) follows. Conversely assume (iii), and write \(N:=\mathrm{Val}(a)=\mathrm{Val}(b)\). Proposition \ref{prop:finite-local-fold-normalization} gives finite symmetric rewrite sequences
\[
a\rightsquigarrow \Fold_\infty(a)=Z(N),
\qquad
b\rightsquigarrow \Fold_\infty(b)=Z(N).
\]
Since the moves are symmetric, the second sequence may be reversed. Concatenating the first sequence with the reverse of the second gives a finite symmetric rewrite sequence from \(a\) to \(b\). Hence (iii) implies (i).

For the final assertion, Proposition \ref{prop:finite-local-fold-normalization} shows that every orbit contains \(\Fold_\infty(a)\). If \(c,c'\in X_\infty\) lie in the same orbit, then their values are equal by the already proved implication (i)\(\Rightarrow\)(iii). The uniqueness part of Theorem \ref{thm:zeckendorf-unique} then gives \(c=c'\). Thus the Zeckendorf representative in the orbit is unique.
\end{proof}

\begin{corollary}[Folding as a computable section of the quotient map]\label{cor:fold-as-section}
The preceding theorem may be summarized as
\[
\mathcal{D}\twoheadrightarrow \mathcal{D}/\!\equiv\ \cong\ \NN
\qquad\text{and}\qquad
X_\infty\hookrightarrow \mathcal{D},
\]
where the inclusion of \(X_\infty\) supplies a computable section selecting the canonical representative in each class. This is the Fibonacci instance of the abstract mechanism
\[
\text{folding}=\text{passing to an invariant/quotient}+\text{choosing normal form}.
\]
The finite local reachability realization of this section is Proposition \ref{prop:finite-local-fold-normalization}.
\end{corollary}

\begin{proof}
This is an immediate restatement of Theorem \ref{thm:monoid-quotient-is-N}, with computability and local finite realization supplied by Proposition \ref{prop:finite-local-fold-normalization}.
\end{proof}
